\section{Liturgical Theology: One Prayer, Whole Pasch}

The General Instruction's account of Eucharistic Prayer keeps the doctrine
from becoming an
abstraction. Thanksgiving addresses the Father for creation and salvation.
The Sanctus joins earthly praise to heavenly worship. The institution
narrative speaks Christ's words within prayer. Anamnesis remembers not an
isolated death but, according to the prayer, Passion, Resurrection, Ascension,
and awaited coming. The oblation, distinctly, offers the one sacrifice under
sacramental signs.
Epiclesis and communion petitions ask that recipients become one body in the
Spirit. Intercessions stretch across living and dead, local and universal
Church. Doxology returns everything to the Father through the Son in the
Spirit (GIRM 79).

The 2002 typical edition of the Roman Missal supplies one concrete instance.
The Roman Canon's \emph{Unde et memores} remembers Passion, Resurrection, and
Ascension before offering the holy victim and gifts. This bounded Roman
evidence is not a universal causal anatomy of consecration. It nonetheless
shows why ``back to Calvary'' is too narrow: the Church celebrates and
sacramentally participates in the Paschal Mystery in praise of the Father and
expectation of fulfillment.

The distinction between representation and reenactment is equally important.
The liturgy does not stage another Last Supper or crucifixion. The altar is not
a theatrical set, the priest is not acting the role of a deceased founder, and
the fraction does not mime Christ being killed again. Sacramental signs effect
what Christ promises in the mode he instituted. Their restraint belongs to
their truth.

